\subsection{交换 Banach 代数的态射}

命题 9. --- 设 A 是一个 Banach 代数，并设 B 是一个交换半单 Banach 代数。从 A 的基础代数到 B 的基础代数的每个态射都是连续的。

设 \(h \colon \mathrm{A} \to \mathrm{B}\) 是一个代数同态，并设 \((a, b) \in \mathrm{A} \times \mathrm{B}\) 是附着于图像 \(\Gamma\)（\(h\) 的图像）的一个点。设 \(\chi \in \mathsf{X}'(\mathrm{B})\)。函数 \(x \mapsto \chi(h(x))\)（从 \(\mathrm{A}\) 到 \(\mathbf{C}\)）是一个代数同态，因而是连续的 (I, p.~29, 定理 1)。于是从 \(\mathrm{A} \times \mathrm{B}\) 到 \(\mathbf{C}\) 中的、由 \((x, y) \mapsto \chi(h(x)) - \chi(y)\) 给出的映射是连续的，它在 \(\Gamma\) 上为零，因而在 \((a, b)\) 处为零。这样，我们有 \(\chi(h(a)) = \chi(b)\) 对一切 \(\chi \in \mathsf{X}'(\mathrm{B})\) 成立。由于 B 是半单的，我们有 \(h(a) = b\)。因而 \(h\) 的图像是闭的，从而 \(h\) 是连续的 (EVT, I, p.~19, 推论 5)。

推论. --- 在一个无根的复交换代数上，定义 Banach 代数结构的两个范数是等价的。

只需对该代数的恒等映射应用命题 9。

设 A 与 B 是交换 Banach 代数。依 I, p.~9, 第 7 小节，若 \(h\colon A\to B\) 是一个满射态射，则 \(\mathsf{X}'(h)\) 是从 \(\mathsf{X}'(\mathsf{B})\) 到 \(\mathsf{X}'(\mathsf{A})\) 的一个闭子空间上的同胚，且它把 0 映到 0。（在 \(h\) 是一个子代数到 \(\mathcal{C}(\mathsf{X})\) 中的单射的情形，在弱得多的假设下映射 \(\mathsf{X}'(h)\) 就是单射，参见 I, p.~142, 命题 1, d)。）

现在设 \(h \colon \mathrm{A} \to \mathrm{B}\) 是一个单射态射。一般说来，\(\mathsf{X}'(h)\) 不是满射的，但下面的命题给出了使之成立的一个必要条件。

命题 10. --- 设 A 与 B 是幺交换 Banach 代数，\(h\colon A\to B\) 是一个保单位元的代数同态，不必连续。若 \(\mathsf{X}(h)\) 是满射的，则 \(h(\mathrm{A})\) 是 B 的一个全子代数。

设 \(x \in A\) 使得 \(h(x)\) 在 B 中可逆。对每个 \(\chi \in \mathsf{X}(\mathsf{A})\)，存在 \(\xi \in \mathsf{X}(\mathsf{B})\) 使得 \(\chi = \mathsf{X}(h)(\xi)\)，因而 \(\chi(x) = \xi(h(x)) \neq 0\)。于是 I, p.~37, 命题 6 表明 x 在 A 中可逆，从而 \(h(x)\) 在 \(h(\mathsf{A})\) 中可逆。

该命题中的必要条件不是充分的，即使 h 是等距的也是如此 (I, p.~168, 习题 14)。不过我们有如下的结果：

命题 11. --- 设 A 与 B 是交换幺 Banach 代数，a 是 A 的一个元素，而 h: A → B 是一个单的保单位元态射（不必连续）。假定由 a 生成的 A 的闭全子代数等于 A。下列条件等价：

\begin{enumerate}
\def\labelenumi{(\roman{enumi})}
\item
  \(\mathsf{X}(h)\) 是满射的；
\item
  \(h(\mathrm{A})\) 是 B 的一个全子代数；
\item
  \(\mathrm{Sp}_{\mathrm{A}}(a)=\mathrm{Sp}_{\mathrm{B}}(h(a)).\)
\item
  (i) \(\Longrightarrow\) (ii) 由命题 10 得出。
\item
  (ii) \(\Longrightarrow\) (iii) 由公式 \(\mathrm{Sp}_{\mathrm{A}}(a) = \mathrm{Sp}_{h(\mathrm{A})}(h(a)) = \mathrm{Sp}_{\mathrm{B}}(h(a))\) 得出，该公式有效是因为 \(h(\mathrm{A})\) 是 B 的一个全子代数。
\item
  (iii) \(\Longrightarrow\) (i) 依 I, p.~6, 第 6 小节的公式 (4)，我们有交换图表
\end{enumerate}

\[
\begin{array}{c} \mathsf {X} (\mathrm{B}) \xrightarrow {\mathsf {X} (h)} \mathsf {X} (\mathrm{A}) \\ \Big \downarrow^ {\mathcal {G} _ {\mathrm{B}} (h (a))} \qquad \qquad \qquad \Big \downarrow^ {\mathcal {G} _ {\mathrm{A}} (a)} \\ \operatorname{Sp} _ {\mathrm{B}} (h (a)) \xrightarrow {i} \operatorname{Sp} _ {\mathrm{A}} (a) \end{array}
\]

其中竖直箭头表示满射 (I, p.~37, 命题 6)，而 \(i\) 是典范包含。假设 (iii) 意味着 \(i\) 是双射。此外，满射 \(\mathcal{G}_{\mathrm{A}}(a)\colon \mathsf{X}(\mathrm{A})\to \mathrm{Sp}_{\mathrm{A}}(a)\) 是双射：事实上，对 A 的任意特征标 \(\chi_{1}\) 与 \(\chi_{2}\)，集合 \(\mathrm{A}_{\chi_1,\chi_2} = \{x\in \mathrm{A}\mid \chi_1(x) = \chi_2(x)\}\) 是 A 的一个闭全子代数。于是依假设，我们有 \(\mathrm{A}_{\chi_1,\chi_2} = \mathrm{A}\)（若 \(a\in \mathrm{A}_{\chi_1,\chi_2}\)），也就是说，\(\chi_{1} = \chi_{2}\)（若 \(\mathcal{G}_{\mathrm{A}}(a)\chi_{1} = \mathcal{G}_{\mathrm{A}}(a)\chi_{2}\)）。于是该图表蕴涵映射 \(\mathsf{X}(h)\) 是满射的。

